% Part: first-order-logic
% Chapter: sequent-calculus
% Section: provability-consistency

\documentclass[../../../include/open-logic-section]{subfiles}

\begin{document}

\iftag{FOL}
      {\olfileid{fol}{seq}{prv}}
      {\olfileid{pl}{seq}{prv}}

\olsection{\usetoken{S}{derivability}, అవైరుధ్యం}

ఇప్పుడు \tetoken{వ్యుత్పాద్యత}{!!{derivability}} సంబంధపు అనేక ధర్మాలను స్థాపిస్తాం. అవి
స్వతంత్రంగా ఆసక్తికరమైనవే కాక, సంపూర్ణతా సిద్ధాంత నిరూపణలో ఒక్కొక్కటి
ఒక పాత్ర పోషిస్తాయి.

\begin{prop}\ollabel{prop:provability-contr}
  $\Gamma \Proves !A$ అయి, $\Gamma \cup \{!A\}$ వైరుధ్యభరితమైతే,
  $\Gamma$ వైరుధ్యభరితం.
\end{prop}

\begin{proof}
$\Gamma_0$, $\Gamma_1 \subseteq \Gamma$ అనే పరిమిత సమితులు ఉండి,
$\Log{LK}$ అనేది $\Gamma_0 \Sequent !A$, $!A, \Gamma_1
\Sequent \quad$లను \tetoken{నిగమనం}{!!{derive}s} చేస్తుంది. $\Gamma_0 \Sequent !A$
యొక్క $\Log{LK}$-\tetoken{వ్యుత్పత్తిని}{!!{derivation}} $\pi_0$ అనీ, $\Gamma_1, !A
\Sequent \quad$ యొక్క $\Log{LK}$-\tetoken{వ్యుత్పత్తిని}{!!{derivation}} $\pi_1$ అనీ
అందాం. అప్పుడు కిందివిధంగా \tetoken{నిగమనం}{!!{derive}} చేయవచ్చు:
\begin{prooftree}
\AxiomC{}
\RightLabel{$\pi_0$}
\Deduce$ \Gamma_0 \fCenter !A $
\AxiomC{}
\RightLabel{$\pi_1$}
\Deduce$!A, \Gamma_1 \fCenter $
\RightLabel{\Cut}
\BinaryInf$ \Gamma_0,\Gamma_1 \fCenter $
\end{prooftree}

$\Gamma_0 \subseteq \Gamma$, $\Gamma_1 \subseteq \Gamma$ కనుక
$\Gamma_0 \cup \Gamma_1 \subseteq \Gamma$; అందువల్ల $\Gamma$
వైరుధ్యభరితం.
\end{proof}

\begin{prop}
\ollabel{prop:prov-incons}
$\Gamma \cup \{\lnot !A\}$ వైరుధ్యభరితమైనప్పుడు, అప్పుడు మాత్రమే
$\Gamma \Proves !A$.
\end{prop}

\begin{proof}
మొదట $\Gamma \Proves !A$ అనుకుందాం. అంటే $\Gamma \Sequent !A$
యొక్క ఒక \tetoken{వ్యుత్పత్తి}{!!a{derivation}}~$\pi_0$ ఉంది. ఒక $\LeftR{\lnot}$ నియమాన్ని
చేర్చితే $\lnot !A, \Gamma \Sequent \quad$ యొక్క ఒక
\tetoken{వ్యుత్పత్తి}{!!a{derivation}} వస్తుంది; అంటే $\Gamma \cup \{\lnot !A\}$
వైరుధ్యభరితం.

$\Gamma \cup \{\lnot !A\}$ వైరుధ్యభరితమైతే,
$\lnot !A, \Gamma \Sequent \quad$ యొక్క ఒక \tetoken{వ్యుత్పత్తి}{!!a{derivation}}~$\pi_1$
ఉంటుంది. కిందిది $\Gamma \Sequent !A$ యొక్క ఒక \tetoken{వ్యుత్పత్తి}{!!a{derivation}}:
\begin{prooftree}
  \Axiom$!A \fCenter !A$
  \RightLabel{\RightR{\lnot}}
  \UnaryInf$\fCenter !A, \lnot !A$
  \AxiomC{}
  \RightLabel{$\pi_1$}
  \Deduce$\lnot !A, \Gamma \fCenter$
  \RightLabel{\Cut}
  \BinaryInf$\Gamma \fCenter !A$
\end{prooftree}
\end{proof}

\begin{prob}
$\Gamma \cup \{!A\}$ వైరుధ్యభరితమైనప్పుడు, అప్పుడు మాత్రమే
$\Gamma \Proves \lnot !A$ అని నిరూపించండి.
\end{prob}

\begin{prop}\ollabel{prop:explicit-inc}
  $\Gamma \Proves !A$, $\lnot !A \in \Gamma$ అయితే $\Gamma$
  వైరుధ్యభరితం.
\end{prop}

\begin{proof}
  $\Gamma \Proves !A$, $\lnot !A \in \Gamma$ అనుకుందాం. అప్పుడు
  $\Gamma_0 \Sequent !A$ అనే ఒక సీక్వెంట్ యొక్క
  \tetoken{ఒక వ్యుత్పత్తి}{!!a{derivation}}~$\pi$ ఉంటుంది. $\lnot !A, \Gamma_0 \Sequent \quad$
  సీక్వెంట్ కూడా \tetoken{వ్యుత్పాదించదగిన}{!!{derivable}}:
  \begin{prooftree}
    \AxiomC{}
    \RightLabel{$\pi$}
    \Deduce$\Gamma_0 \fCenter !A$
    \Axiom$!A \fCenter !A$
    \RightLabel{\LeftR{\lnot}}
    \UnaryInf$\lnot !A, !A \fCenter$
    \RightLabel{\LeftR{\Exchange}}
    \UnaryInf$!A, \lnot !A \fCenter$
    \RightLabel{\Cut}
    \BinaryInf$\Gamma_0, \lnot !A \fCenter$
  \end{prooftree}
  $\lnot !A \in \Gamma$, $\Gamma_0 \subseteq \Gamma$ కనుక ఇది
  $\Gamma$ వైరుధ్యభరితమని చూపుతుంది.
\end{proof}

\begin{prop}\ollabel{prop:provability-exhaustive}
  $\Gamma \cup \{!A\}$, $\Gamma \cup \{\lnot !A\}$ రెండూ
  వైరుధ్యభరితమైతే, $\Gamma$ వైరుధ్యభరితం.
\end{prop}

\begin{proof}
$\Gamma_0 \subseteq \Gamma$, $\Gamma_1 \subseteq \Gamma$ అనే పరిమిత
సమితులు, వరుసగా $!A, \Gamma_0 \Sequent \quad$,
$\lnot !A, \Gamma_1 \Sequent \quad$ల యొక్క $\Log{LK}$-\tetoken{వ్యుత్పత్తులు}{!!{derivation}s}
$\pi_0$, $\pi_1$ ఉన్నాయి. అప్పుడు కిందివిధంగా \tetoken{నిగమనం}{!!{derive}} చేయవచ్చు:
\begin{prooftree}
\AxiomC{}
\RightLabel{$\pi_0$}
\Deduce$ !A, \Gamma_0 \fCenter $
\RightLabel{\RightR{\lnot}}
\UnaryInf$ \Gamma_0 \fCenter \lnot !A$
\AxiomC{}
\RightLabel{$\pi_1$}
\Deduce$\lnot !A, \Gamma_1 \fCenter  $
\RightLabel{\Cut}
\BinaryInf$ \Gamma_0, \Gamma_1 \fCenter $
\end{prooftree}
$\Gamma_0 \subseteq \Gamma$, $\Gamma_1 \subseteq \Gamma$ కనుక
$\Gamma_0 \cup \Gamma_1 \subseteq \Gamma$. అందువల్ల $\Gamma$
వైరుధ్యభరితం.
\end{proof}

\end{document}
